\documentclass{zhvt-classic}

% 使用本地源流明体（GenRyuMin2 TW）按文件路径加载，不依赖系统字体名
% 版框（内框／界栏／书口侧栏／鱼尾）改用与版心印章同一印泥棕色：
%   印章图片 yshs.png 实色 = RGB(135,68,52) = HTML 874434
\usepackage{xcolor}
\usepackage{tikz}
% sealbrown 双注册：xcolor 供 TikZ/pgf，L3 供 cls 版框
\definecolor{sealbrown}{HTML}{874434}
\ExplSyntaxOn
\color_set:nnn { sealbrown } { HTML } { 874434 }
\ExplSyntaxOff

\zhvtset{
  grid color       = sealbrown,
  main font        = GenRyuMin2TW-SB,
  main font path   = C:/Users/Administrator/texmf/fonts/opentype/public/GenRyuMin2TW/,
jiazhu font      = GenRyuMin2TW-M,
  jiazhu font path = C:/Users/Administrator/texmf/fonts/opentype/public/GenRyuMin2TW/,
  % 页码框容纳 4 字：正常页「九十七」+前置空 1 格=4 行恰好；
  % 100+ 页「一百零一」共 5 行超出，末字「一/二/三」被框裁掉（2026-10-04 实测），故放宽到 5
  shukou page number chars = 5,
  % 书口书名框高从默认 6 字减到 4 字，把省出的 2 字高拨给 chapter 框，
  % 使 6 字书名+卷名放进 9 字高的 chapter 框后不再压线（2026-10-05）
  shukou title chars = 4,
}

\title{史记}[魏公子列传]
\author{汉·司马迁}
\maker{闲叟漫记·重製}


% ===== 索隐述赞样式（2026-11-08 四次定稿 · 愚大人亲手改定参数）=====
% ⚠ 竖排前提：自定义字体家族必须带 RawFeature={vertical:+vert}；
%   simhei.ttf 加 +vhal 会整体加载失败静默回退，只写 +vert；simfang 同理。
\newCJKfontfamily{\shuzanfont}[Path=C:/Windows/Fonts/,RawFeature={vertical:+vert}]{simhei.ttf}
\newCJKfontfamily{\shuzanbody}[Path=C:/Windows/Fonts/,RawFeature={vertical:+vert}]{simfang.ttf}
% 题头：152x30pt 圆角 6pt 棕底白字；baseline (0,2) 使整框落入界行格子不压版框线；
% 白字节点 x 不补偿、y +8.2pt（像素级墨迹实测 dx/dy≈0，L/R 留白约 20pt）。
\newcommand{\shuzanhead}[1]{%
  \begin{tikzpicture}[baseline={(0,2)}]
    \fill[sealbrown, rounded corners=6pt]
      (-76pt,-15pt) rectangle (76pt,15pt);
    \node[text=white, inner sep=0pt, anchor=center] at (0pt,8.2pt)
      {\shuzanfont\fontsize{23}{23}\selectfont #1};
  \end{tikzpicture}%
}
% ★ 修正书口卷名位置：chapter 框起始位置改回 cls 默认值，避免全局影响。
%   只在每个卷首页通过 \juan 命令里的空命令微调该页位置（安全，不会跨页污染）。
\ExplSyntaxOn
\AtBeginDocument{
  \dim_set:Nn \l__zhvt_shukou_chapter_title_start_dim
    { \l__zhvt_shukou_title_height_dim + \l__zhvt_font_size_dim * 2 }
  \dim_set:Nn \l__zhvt_shukou_page_number_start_dim
    { \l__zhvt_shukou_chapter_title_start_dim + \l__zhvt_shukou_chapter_title_height_dim }
}
\ExplSyntaxOff

\begin{document}

% ★★ 版心（书口）排布覆写：把书名放到鱼尾之下（2026-10-04）
%
%   cls 默认版心自上而下：书名（title 静态框，在鱼尾之上）→ 鱼尾 → 篇名（chapter 动态框）→ 页码
%   本模板改为：          鱼尾（位置不动）→ 书名 → 篇名 → 页码分割线 → 页码
%
%   手法：把 cls 放在 title 静态框里的书名，改放进 chapter 动态框（鱼尾之下）；
%         title 静态框留空。书名取 \l__zhvt_title_str，即 \title{...} 的第一参数，
%         所以换书时只改 \title，不必在此重复写书名。
%
%   ⚠ 三条必须遵守的写法（2026-10-04 实测）：
%     ① 汉字必须用 \hbox{...} 包住。flowfram 会把框内容装进一个宽 14.4mm（= 一个 line_sep
%        = 书口栏宽）的纵向 minipage 并按该宽度折行，而整页靠 PDF /Rotate 90 显示、
%        框内文字本身不变向 —— 不包 \hbox 时「围炉夜话」4 字会被逐字折成 4 行，
%        在版面上就成了横排 4 字并横向溢出版框。
%     ② 竖排下**一个 \hbox 就是一整列**。书名与篇名要同处一列，必须放进同一个 \hbox；
%        分成两个 \hbox 时，后一个会被当成新的一列、甩到隔壁栏去
%        （实测「序」跑到书口栏外侧 14.4mm 的正文列上）。中间的 \quad 在框内是横向空白，
%        显示出来正好是书名与篇名之间的纵向间隔。
%     ③ ★ 内容盒要写成 **与框等宽的定宽盒** `\hbox to \l__zhvt_shukou_chapter_title_height_dim{…\hfil}`。
%        实测：chapter 框是把内容盒沿纵向**居中**摆放的，内容宽度一变，书名就跟着漂：
%        自然宽度的 \hbox（书名 38.4mm）落在 x=157.2mm，加篇名后（57.6mm）落在 x=152.4mm，
%        目录页（无篇名）与正文页的书名就差半格。做成等宽盒后，无论有没有篇名，
%        内容固定从框顶起排，实测两页均落在 x=142.8mm，不再漂。
\ExplSyntaxOn
% 把书口 chapter 框（书名+篇名）起始位置下移 1 字高，避免 4 字卷名顶到鱼尾
% ⚠ cls 的 \__zhvt_calculate_dimensions: 在 \AtBeginDocument 里执行，会重算该 dim；
%   必须把覆写挂到 begin document 之后，否则会被二次覆盖（2026-10-05 实测）。
% (旧方案已移到导言区，直接修改 chapter 框起始位置)
% 卷标题命令 #1=书口篇名(鱼尾) #2=正文卷名（含卷序，卷名内空格用 ~）
%
% ★★ 目录页两项修订（2026-10-05，愚大人指出）★★
%   ① 页码起点统一对齐：
%      原状 —— 目录项写成 \par\hspace{4em}卷名\hspace{4em}页码，页码起点 = 卷名实宽 + 4em。
%      卷名实宽随字数变：「卷一~梁惠王上」6 字 +~、「卷七~离娄上」5 字 +~，
%      差 1 个字位 → 卷七~卷十 4 条的页码整体比其余 10 条高 27.21pt（一个字），看着上下不齐。
%      改法 —— 卷名补全角空格（宽 = 1 字 = 27.215pt）补足到 **7 字符位**再写进目录，
%      14 条页码起点遂全部落在同一条竖线上（起点对齐；末字随页码字数自然参差）。
%      ⚠ 只对目录里的副本 \l_tmpc_tl 补，正文卷名仍用 \l_tmpa_tl 原样 —— 正文像素不受影响。
%   ② 卷名整体上移半字位：前导 \hspace 由 4em 改 3.5em。
%      实测原状目录块在版框内偏下（天头留白 139.6pt / 地脚留白 110.7pt），
%      上移 0.5em = 13.6pt 后上下基本均分，视觉居中。
\NewDocumentCommand { \juan } { mm } {
	\clearpage
	\tl_set:Nn \l__zhvt_chapter_title_tl { #1 }
	\tl_set:Nn \l_tmpa_tl { #2 }
	\tl_set_eq:NN \l_tmpc_tl \l_tmpa_tl
	\int_step_inline:nn { 7 - \str_count:n { #2 } }
	  { \tl_put_right:Nn \l_tmpc_tl { 　 } }
	\tl_set:Nx \l_tmpb_tl { \zhnum { page } }
	\tl_put_left:Nn \l_tmpb_tl { \hspace { 4em } }
	\tl_put_left:Nn \l_tmpb_tl { \l_tmpc_tl }
	\tl_put_left:Nn \l_tmpb_tl { \par \hspace { 3.5em } }
	\addtocontents { toc } { \l_tmpb_tl }
	\skip_horizontal:n { \l__zhvt_font_size_dim * 4 }
	\l_tmpa_tl
	\par
	\thispagestyle { empty }
}
\cs_set_protected:Npn \__zhvt_set_shukou_content: {
	% ★ 鱼尾上方留白（title 静态框）嵌入闲叟钤印：
	%   整页 /Rotate 90，故此处把图用 \rotatebox{90} 转正为竖章；
	%   尺寸卡在版心条幅内（书内约 7mm 宽 × 33.8mm 高，留余量不压鱼尾）。
	\setstaticcontents * { title } {
		\vbox_to_ht:nn { \l__zhvt_baseline_skip_dim } {
			\vskip 14.2pt
			\hbox_to_wd:nn { \l__zhvt_shukou_title_height_dim } {
				\hfil
				\rotatebox [ origin = c ] { 90 } { \includegraphics [ width = 7mm , height = 33.8mm ] { yshs.png } }
				\hfil
			}
		}
	}
% ★ 版心（书口）统一改法（2026-11-16 愚大人定）：18pt＋框内居中＋书口文字同版框色
%   ① 原 27.2pt 时 7 字长篇名（如「廉颇蔺相如列传」）末字溢出 chapter 框、压住页码；
%      改 \fontsize{18pt}{18pt}，前后 \hfil 令字串在框内居中（不随篇名字数漂移）。
%   ② \vbox{\vskip 8.2pt} 定量下移，抵消缩字号后栏宽方向 8.2pt 上偏（像素实测 y 中心 543.3 = 原版）。
%      ⚠ \raisebox / \vspace* 在 flowfram 框内无效；换字号须同步换补偿：16pt→10.1、19pt→7.4。
%   ③ chapter 与 page~number 均置 \color{sealbrown}（HTML 874434），与版框／界行同色；正文仍黑。
%   ⚠ 页码字位、钤印、鱼尾、界行、正文一律未动；五部合订本同此口径。
	\setdynamiccontents * { chapter }
	{ \color{sealbrown} \vbox { \vskip 8.2pt \hbox to \l__zhvt_shukou_chapter_title_height_dim
		{ \hfil \fontsize{18pt}{18pt}\selectfont \l__zhvt_title_str \quad \l__zhvt_chapter_title_tl \hfil } } }
	\setdynamiccontents * { page~number } { \color{sealbrown} \quad \zhnum { page } }
}
\__zhvt_set_shukou_content:
\ExplSyntaxOff

\maketitle

% 单篇，无多卷目录，不渲染目录页

\juan{魏公子列传}{卷七十七~魏公子列传第十七}
\hspace*{2em}魏公子无忌者。魏昭王少子而魏安釐王异母弟也。昭王薨。安釐王即位。封公子为信陵君。是时范睢亡魏相秦。以怨魏齐故。秦兵围大梁。破魏华阳下军。走芒卯。魏王及公子患之。\\
\hspace*{2em}公子为人仁而下士。士无贤不肖皆谦而礼交之。不敢以其富贵骄士。士以此方数千里争往归之。致食客三千人。当是时。诸侯以公子贤。多客。不敢加兵谋魏十馀年。\\
\hspace*{2em}公子与魏王博。而北境传举烽。言﹁赵寇至。且入界﹂。魏王释博。欲召大臣谋。公子止王曰。﹁赵王田猎耳。非为寇也。﹂复博如故。王恐。心不在博。居顷。复从北方来传言曰。﹁赵王猎耳。非为寇也。﹂魏王大惊。曰。﹁公子何以知之。﹂公子曰。﹁臣之客有能深得赵王阴事者。赵王所为。客辄以报臣。臣以此知之。﹂是後魏王畏公子之贤能。不敢任公子以国政。\\
\hspace*{2em}魏有隐士曰侯嬴。年七十。家贫。为大梁夷门监者。公子闻之。往请。欲厚遗之。不肯受。曰。﹁臣脩身絜行数十年。终不以监门困故而受公子财。﹂公子於是乃置酒大会宾客。坐定。公子从车骑。虚左。自迎夷门侯生。侯生摄敝衣冠。直上载公子上坐。不让。欲以观公子。公子执辔愈恭。侯生又谓公子曰。﹁臣有客在市屠中。原枉车骑过之。﹂公子引车入巿。侯生下见其客朱亥。俾倪故久立。与其客语。微察公子。公子颜色愈和。当是时。魏将相宗室宾客满堂。待公子举酒。巿人皆观公子执辔。从骑皆窃骂侯生。侯生视公子色终不变。乃谢客就车。至家。公子引侯生坐上坐。遍赞宾客。宾客皆惊。酒酣。公子起。为寿侯生前。侯生因谓公子曰。﹁今日嬴之为公子亦足矣。嬴乃夷门抱关者也。而公子亲枉车骑。自迎嬴於众人广坐之中。不宜有所过。今公子故过之。然嬴欲就公子之名。故久立公子车骑巿中。过客以观公子。公子愈恭。巿人皆以嬴为小人。而以公子为长者能下士也。﹂於是罢酒。侯生遂为上客。\\
\hspace*{2em}侯生谓公子曰。﹁臣所过屠者硃亥。此子贤者。世莫能知。故隐屠间耳。﹂公子往数请之。硃亥故不复谢。公子怪之。\\
\hspace*{2em}魏安釐王二十年。秦昭王已破赵长平军。又进兵围邯郸。公子姊为赵惠文王弟平原君夫人。数遗魏王及公子书。请救於魏。魏王使将军晋鄙将十万众救赵。秦王使使者告魏王曰。﹁吾攻赵旦暮且下。而诸侯敢救者。已拔赵。必移兵先击之。﹂魏王恐。使人止晋鄙。留军壁鄴。名为救赵。实持两端以观望。平原君使者冠盖相属於魏。让魏公子曰。﹁胜所以自附为婚姻者。以公子之高义。为能急人之困。今邯郸旦暮降秦而魏救不至。安在公子能急人之困也。且公子纵轻胜。弃之降秦。独不怜公子姊邪。﹂公子患之。数请魏王。及宾客辩士说王万端。魏王畏秦。终不听公子。公子自度终不能得之於王。计不独生而令赵亡。乃请宾客。约车骑百馀乘。欲以客往赴秦军。与赵俱死。\\
\hspace*{2em}行过夷门。见侯生。具告所以欲死秦军状。辞决而行。侯生曰。﹁公子勉之矣。老臣不能从。﹂公子行数里。心不快。曰。﹁吾所以待侯生者备矣。天下莫不闻。今吾且死而侯生曾无一言半辞送我。我岂有所失哉。﹂复引车还。问侯生。侯生笑曰。﹁臣固知公子之还也。﹂曰。﹁公子喜士。名闻天下。今有难。无他端而欲赴秦军。譬若以肉投馁虎。何功之有哉。尚安事客。然公子遇臣厚。公子往而臣不送。以是知公子恨之复返也。﹂公子再拜。因问。侯生乃屏人间语。曰。﹁嬴闻晋鄙之兵符常在王卧内。而如姬最幸。出入王卧内。力能窃之。嬴闻如姬父为人所杀。如姬资之三年。自王以下欲求报其父仇。莫能得。如姬为公子泣。公子使客斩其仇头。敬进如姬。如姬之欲为公子死。无所辞。顾未有路耳。公子诚一开口请如姬。如姬必许诺。则得虎符夺晋鄙军。北救赵而西卻秦。此五霸之伐也。﹂公子从其计。请如姬。如姬果盗晋鄙兵符与公子。\\
\hspace*{2em}公子行。侯生曰。﹁将在外。主令有所不受。以便国家。公子即合符。而晋鄙不授公子兵而复请之。事必危矣。臣客屠者硃亥可与俱。此人力士。晋鄙听。大善。不听。可使击之。﹂於是公子泣。侯生曰。﹁公子畏死邪。何泣也。﹂公子曰。﹁晋鄙嚄唶宿将。往恐不听。必当杀之。是以泣耳。岂畏死哉。﹂於是公子请硃亥。硃亥笑曰。﹁臣乃市井鼓刀屠者。而公子亲数存之。所以不报谢者。以为小礼无所用。今公子有急。此乃臣效命之秋也。﹂遂与公子俱。公子过谢侯生。侯生曰。﹁臣宜从。老不能。请数公子行日。以至晋鄙军之日。北乡自刭。以送公子。﹂公子遂行。\\
\hspace*{2em}至鄴。矫魏王令代晋鄙。晋鄙合符。疑之。举手视公子曰。﹁今吾拥十万之众。屯於境上。国之重任。今单车来代之。何如哉。﹂欲无听。硃亥袖四十斤铁椎。椎杀晋鄙。公子遂将晋鄙军。勒兵下令军中曰。﹁父子俱在军中。父归。兄弟俱在军中。兄归。独子无兄弟。归养。﹂得选兵八万人。进兵击秦军。秦军解去。遂救邯郸。存赵。赵王及平原君自迎公子於界。平原君负籣矢为公子先引。赵王再拜曰。﹁自古贤人未有及公子者也。﹂当此之时。平原君不敢自比於人。公子与侯生决。至军。侯生果北乡自刭。\\
\hspace*{2em}魏王怒公子之盗其兵符。矫杀晋鄙。公子亦自知也。已卻秦存赵。使将将其军归魏。而公子独与客留赵。赵孝成王德公子之矫夺晋鄙兵而存赵。乃与平原君计。以五城封公子。公子闻之。意骄矜而有自功之色。客有说公子曰。﹁物有不可忘。或有不可不忘。夫人有德於公子。公子不可忘也。公子有德於人。原公子忘之也。且矫魏王令。夺晋鄙兵以救赵。於赵则有功矣。於魏则未为忠臣也。公子乃自骄而功之。窃为公子不取也。﹂於是公子立自责。似若无所容者。赵王埽除自迎。执主人之礼。引公子就西阶。公子侧行辞让。从东阶上。自言罪过。以负於魏。无功於赵。赵王侍酒至暮。口不忍献五城。以公子退让也。公子竟留赵。赵王以鄗为公子汤沐邑。魏亦复以信陵奉公子。公子留赵。\\
\hspace*{2em}公子闻赵有处士毛公藏於博徒。薛公藏於卖浆家。公子欲见两人。两人自匿不肯见公子。公子闻所在。乃间步往从此两人游。甚欢。平原君闻之。谓其夫人曰。﹁始吾闻夫人弟公子天下无双。今吾闻之。乃妄从博徒卖浆者游。公子妄人耳。﹂夫人以告公子。公子乃谢夫人去。曰。﹁始吾闻平原君贤。故负魏王而救赵。以称平原君。平原君之游。徒豪举耳。不求士也。无忌自在大梁时。常闻此两人贤。至赵。恐不得见。以无忌从之游。尚恐其不我欲也。今平原君乃以为羞。其不足从游。﹂乃装为去。夫人具以语平原君。平原君乃免冠谢。固留公子。平原君门下闻之。半去平原君归公子。天下士复往归公子。公子倾平原君客。\\
\hspace*{2em}公子留赵十年不归。秦闻公子在赵。日夜出兵东伐魏。魏王患之。使使往请公子。公子恐其怒之。乃诫门下。﹁有敢为魏王使通者。死。﹂宾客皆背魏之赵。莫敢劝公子归。毛公。薛公两人往见公子曰。﹁公子所以重於赵。名闻诸侯者。徒以有魏也。今秦攻魏。魏急而公子不恤。使秦破大梁而夷先王之宗庙。公子当何面目立天下乎。﹂语未及卒。公子立变色。告车趣驾归救魏。\\
\hspace*{2em}魏王见公子。相与泣。而以上将军印授公子。公子遂将。魏安釐王三十年。公子使使遍告诸侯。诸侯闻公子将。各遣将将兵救魏。公子率五国之兵破秦军於河外。走蒙骜。遂乘胜逐秦军至函谷关。抑秦兵。秦兵不敢出。当是时。公子威振天下。诸侯之客进兵法。公子皆名之。故世俗称魏公子兵法。\\
\hspace*{2em}秦王患之。乃行金万斤於魏。求晋鄙客。令毁公子於魏王曰。﹁公子亡在外十年矣。今为魏将。诸侯将皆属。诸侯徒闻魏公子。不闻魏王。公子亦欲因此时定南面而王。诸侯畏公子之威。方欲共立之。﹂秦数使反间。伪贺公子得立为魏王未也。魏王日闻其毁。不能不信。後果使人代公子将。公子自知再以毁废。乃谢病不朝。与宾客为长夜饮。饮醇酒。多近妇女。日夜为乐饮者四岁。竟病酒而卒。其岁。魏安釐王亦薨。\\
\hspace*{2em}秦闻公子死。使蒙骜攻魏。拔二十城。初置东郡。其後秦稍蚕食魏。十八岁而虏魏王。屠大梁。\\
\hspace*{2em}高祖始微少时。数闻公子贤。及即天子位。每过大梁。常祠公子。高祖十二年。从击黥布还。为公子置守冢五家。世世岁以四时奉祠公子。\\
\hspace*{2em}太史公曰。吾过大梁之墟。求问其所谓夷门。夷门者。城之东门也。天下诸公子亦有喜士者矣。然信陵君之接岩穴隐者。不耻下交。有以也。名冠诸侯。不虚耳。高祖每过之而令民奉祠不绝也。\\

% ===== 索隐述赞（司马贞《史记索隐》）· 统一样式 =====
% 题头「索隐述赞曰」五字并入反白框；正文每四句一列，定宽盒 \hfil 居列中；
% \fontsize{24}{41} 列距对齐全书界行（121+41k）；\clearpage 使述赞独占一页。
\clearpage
\shuzanhead{索隐述赞曰}\linebreak
{\shuzanbody\fontsize{24}{41}\selectfont
\hbox to \hsize{\hfil 信陵下士。邻国相倾。以公子故。不敢加兵。\hfil}\linebreak
\hbox to \hsize{\hfil 颇知硃亥。尽礼侯嬴。遂卻晋鄙。终辞赵城。\hfil}\linebreak
\hbox to \hsize{\hfil 毛。薛见重。万古希声。\hfil}%
\par}
\end{document}
